\section{\name{}}
\label{sec:klear}
%\begin{figure*}[t]
\centering
\includegraphics[width=\textwidth]{Figures/raw_figures/figure1.pdf}
%\vspace{-0.2in}
\caption{
\textbf{Overview of \name{}.}
\textbf{(a)} Seed questions from Korean benchmarks are labeled by knowledge (K) and application (A) demands.
\textbf{(b)} Silver reasoning is generated from the seed question and gold answer, then used to derive an anchor knowledge bundle for generating higher-A questions.
\textbf{(c)} Generated questions undergo LLM auditing, proxy screening, and human validation, with failed candidates revised using LLM or human feedback.
The result is a matched set of questions that share the same knowledge but differ in application demand.
}\label{fig:overview}
\vspace{-0.1in}
\end{figure*}

In this section, we propose \textbf{\name{}} (\textbf{K}nowledge-Matched \textbf{L}LM \textbf{E}valuation in Korean across \textbf{A}pplication \textbf{R}equirements), a framework that constructs knowledge-matched contrasts from existing benchmark questions in three stages (Figure~\ref{fig:overview}).
First, anchor-based generation (Sec.~\ref{sec:generation}) makes the knowledge behind a source question explicit as an \textit{anchor bundle}, and generates new questions that rest on the same bundle while demanding more application.
Second, automated audit and revision (Sec.~\ref{sec:audit}) uses rubric-based LLM evaluation to check the validity and quality of each candidate, screens its difficulty with a proxy model, and revises failing candidates.
Third, human validation and revision (Sec.~\ref{sec:human}) makes a final check of its quality, and revises rejected questions with annotator rationales.
By fixing the knowledge before generation and checking its preservation afterwards, KLEAR isolates differences in application demand between source questions and their contrasts.

\subsection{Anchor-Based Generation}
\label{sec:generation}

\paragraph{Seed selection.}
We draw seeds from the questions profiled in Sec.~\ref{sec:motivation}.
We select A1 or A2 questions that at least two of the three Korean LLMs answer incorrectly, so that construction starts from questions these models find difficult.
An A1 seed can target A2 or A3, whereas an A2 seed can target A3.

\paragraph{Anchor bundle.}
For each seed, an LLM constructs an \textit{anchor bundle} using the question, answer options, gold answer, and silver reasoning.
The bundle consists of three components: 
(1) \textit{Knowledge units} specify the rules, criteria, definitions, or facts required to solve the question, including those underlying the distractors.
(2) \textit{Relations} describe how these units interact, such as which unit is the principle and which the exception.
(3) A \textit{Scope note} explicitly identifies knowledge excluded from the bundle.
Together, these components define a closed set of domain knowledge that subsequent generation must stay within.
The LLM also assesses whether the bundle can support each target-application level, excluding unsupported targets from generation.
For example, an A3 target requires at least two knowledge units that could plausibly govern the same case.
% We call a generated question \textit{knowledge-matched} to its seed when both can be answered from the same anchor bundle and the case alone.

\paragraph{Target-A generation.}
Given a seed, its anchor bundle, and a target-application level, the generator LLM  produces a new multiple-choice question with four or five answer options.
The generated question must preserve the answer type of the seed and must not state or paraphrase the anchor knowledge.
For an A2 target, exactly one knowledge unit from the bundle applies, but its applicability must be inferred from the case facts.
For an A3 target, two or more knowledge units from the bundle compete over a single case with comparable supporting evidence.
Each distractor must represent an identifiable error and all options must be parallel in form.
We generate three candidates per seed--target pair.

\subsection{Automated Audit and Revision}
\label{sec:audit}

\paragraph{LLM rubric.}
An LLM auditor receives the seed, its anchor bundle, the target-application level, and the candidate.
It evaluates the candidate against eleven criteria, labeling each as \textsc{pass}, \textsc{fail}, or \textsc{uncertain} with a rationale.
L1--L6 assess question validity, while L7--L11 assess knowledge preservation and other quality criteria.
The full rubric is provided in Table~\ref{tab:rubrics}.

\paragraph{Proxy screening.}
The rubric does not measure difficulty, and a candidate can pass every criterion while being easy to answer.
We therefore use a weaker proxy model to screen for insufficiently challenging candidates.
This proxy answers each candidate from the question and answer options alone, and a candidate it answers correctly is treated as too easy.

\paragraph{Self-revision.}
A candidate is returned to the generator when any criterion is not \textsc{pass} or when the proxy answers it correctly.
If any criterion is not \textsc{pass}, the generator receives every such criterion with its rationale and must resolve them without regressing other criteria.
If all criteria are \textsc{pass} but the proxy answers correctly, the generator is told that a weaker model answered the question correctly and must close the shortcut.
In both cases, the revision must stay within the bundle.
Each candidate is revised at most once, and the revision is audited and screened again under the same conditions.
A candidate passes the automated stage when, in its original or revised form, all eleven criteria are \textsc{pass} and the proxy answers it incorrectly.


\subsection{Human Validation and Revision}
\label{sec:human}

\paragraph{Targeted validation.}
To identify defects that may remain after automated auditing, human annotators evaluate each candidate that passes the automated stage against five criteria in two stages (Table~\ref{tab:rubrics}).
In the first stage, without the anchor bundle, H1--H3 check the surface quality of the question, and in the second stage, with the bundle revealed, H4--H5 check the control of knowledge and application.
Annotators label each criterion as \textsc{pass}, \textsc{borderline}, or \textsc{fail}, with a rationale for any label other than \textsc{pass}.
A candidate passes human validation when every annotator labels all five criteria as \textsc{pass}.

\paragraph{Human-guided revision.}
Candidates that fail to pass are returned to the generator with the human annotators' labels and rationales.
The generator revises each candidate to address the identified defects while remaining within the anchor bundle and preserving the target-application level.
The revised candidates undergo human validation again under the same two-stage protocol.
Only candidates that pass human validation, either initially or after revision, are retained in the final benchmark.